\section{System and Release Protocol}
\label{sec:system}
\label{sec:impl}
\subsection{A Release Crossing Two Availability Boundaries}
\label{sec:motivation}
Our running example has four packages: \code{assets} is a shared npm library,
\code{ui} is an npm package built from that library, \code{base} packages the
UI output into a Docker image, and \code{app} builds from that image. Arrows in Fig.~\ref{fig:stages} run from
provider to consumer. The second builder obtains its base through a registry;
we do not assume that it shares the first builder's local image store.

\begin{figure*}[t]
\centering
\begin{tikzpicture}[>=Latex,node distance=16mm,
 every node/.style={font=\small},
 stage/.style={draw,rounded corners=2pt,minimum width=28mm,minimum height=9mm,align=center}]
\node[stage] (ab) at (0,0) {Build assets\\npm output};
\node[stage] (ub) at (4.3,0) {Build ui\\npm output};
\node[stage] (up) at (4.3,-1.65) {Publish ui};
\node[stage] (bb) at (8.6,0) {Build base\\Docker image};
\node[stage] (cb) at (12.9,0) {Build app\\\texttt{FROM base:v}};
\node[stage] (ap) at (0,-1.65) {Publish assets};
\node[stage] (bp) at (8.6,-1.65) {Publish base};
\node[stage] (cp) at (12.9,-1.65) {Publish app};
\draw[->] (ab) -- node[above,font=\scriptsize]{local output} (ub);
\draw[->] (ub) -- node[above,font=\scriptsize]{local output} (bb);
\draw[->] (ub) -- (up);
\draw[->] (ab) -- (ap);
\draw[->] (bb) -- (bp);
\draw[->,thick] (bp) -- node[above,sloped,font=\scriptsize]{registry required} (cb);
\draw[->] (cb) -- (cp);
\draw[->] (ap) -- node[below,font=\scriptsize]{publish order} (up);
\draw[->] (up) -- node[below,font=\scriptsize]{publish order} (bp);
\draw[->] (bp) -- node[below,font=\scriptsize]{publish order} (cp);
\end{tikzpicture}
\caption{Stage dependencies, not a measured timeline. Publishing \code{assets}
may overlap building \code{ui}, whose publication may overlap building
\code{base}. Building \code{app} must wait for
\code{base}'s publication. Availability modes are configured requirements:
an npm dependency fetched from a registry would also require publication.}
\label{fig:stages}
\end{figure*}

Availability modes describe integration choices, not ecosystems. Here,
\code{ui} and \code{base} use workspace output; \code{app} pulls its base
from a registry. The provider setting \code{isBuildWaitingPublish: true}
therefore belongs on \code{base}. Both npm providers keep the default
local-build mode. Other integrations may require npm publication or share
local Docker images. Consumer publication follows provider publication in
both modes.

Suppose \code{assets} is at 1.4.2, \code{ui} at 3.2.0, \code{base} at
0.8.2, and \code{app} at 2.1.0. An intent to release an assets feature and patch all
transitive consumers plans versions 1.5.0, 3.2.1, 0.8.3, and 2.1.1. After the first
three packages publish and are recorded, \code{app} fails. A planner that only
propagates from \emph{currently unreleased providers} can now lose
\code{app}, leaving an orphan in the release ledger. Our planner instead
asks whether \code{app} has received that intent. Its own pending window is
unchanged, so a retry still owes version 2.1.1.

\subsection{From Repository Inputs to a Plan}
\label{sec:ccme}
The planner reads a fixed repository revision, package definitions, dependency
edges, commit messages, and release tags. Configuration supplies build and
publish commands and availability requirements; manifests can supply package
and dependency identities. The plan specifies selected packages, next
versions, and reasons before any external publication occurs. Changed
inputs require a new plan; a preview is not an authorization for arbitrary
later repository states.

Release intent extends Conventional Commits~\cite{cc}. For example,
\code{feat(assets)\^{}\^{}: add theme assets} directly requests a minor
release of \code{assets} and propagates patch intent to transitive consumers.
An unmarked unit affects its own scope only. A scope can explicitly name
packages or be derived from changed paths. A commit can contain multiple
independent units, each retaining its own source set and propagation depth.
The formal model consumes normalized units, not a particular parser syntax.
The complete language grammar is outside this paper~\cite{ccme2}.
A diagnostics command validates literal intent syntax and parser policy
without committing; it does not verify scope or plan correctness.

Release-worthy descriptions can form version-specific notes, including for
private projects without changelogs. Humans can adapt them for social drafts
or management reports; this unmeasured use does not automate posting.
Notes are narratives; tags discharge release obligations.

Version bumps and prerelease channels are separate axes; a channel transition
need not request a bump. Cancellation, correction, holds, and exact versions
govern pending intent without unpublishing artifacts. These controls motivate
the model predicates; their CLI forms are outside our contribution.

\subsection{Execution, Recording, and Recovery}
For each selected package, the executor schedules version reconciliation,
build, and publish work. A consumer's version/build work waits for provider
builds, or for provider publications when their availability mode requires it.
Build and publish tasks have separate concurrency limits. Independent ready
tasks may proceed concurrently; deterministic planning does not require a
fixed temporal order for independent completion events.

Successful publication creates a per-package completion obligation. The model
records it immediately as a tag. The implementation also has configured
recording and finalization steps, whose failures must be reported; their
crash boundaries are not assumed atomic with an arbitrary publish command.
The next invocation reconstructs outstanding work from the resulting records.
This avoids a separate persistent recovery database, but depends on preserving
and making the release tags available. A lock prevents cooperating release
jobs from writing the same remote release concurrently. It does not make
registry operations transactional or prevent an administrator from rewriting
history outside the protocol.

Two failure cases must remain distinct. If a publish is known to have failed,
its obligation stays pending. If the external operation succeeds but its
response or completion record is lost, an absent tag cannot distinguish that
success from failure. The integration must reconcile the artifact's identity
or use a publisher that safely accepts repetition before recording completion.
Forward recovery does not by itself solve this uncertainty~\cite{helland12}.
The mathematical claims below refer to the recorded state and make their
retry assumptions explicit.

\subsection{Fleet Releases Across Repository Histories}
\label{sec:fleet}
The running chain can span repositories: one owns the npm packages, another
the images. The package graph still orders stages, but a consumer's tag must
also establish which source revisions its release incorporated. The current
language distinguishes three history modes~\cite{ccme2}. Single-history mode
covers one package, a monorepo, or a control-history wrapper. Specified-history
mode takes participating repositories from entry configuration. Discovered-history
mode walks links among peers that own their configurations; a star-shaped
link topology does not make its centre a policy owner.

\paragraph{Ownership and admission.}
Packages build, publish, and record tags in their owning repositories;
central configuration does not relocate publication. Commits are identified by
repository and full object ID, with ancestry confined to that repository.
Cross-repository admission asks whether the consumer has incorporated that
source revision. A control checkpoint associating the exact consumer tag with
source pointers, or a chain of links in the consumer's release record, proves
the boundary. Missing evidence requires an explicit consumer-tag/source-revision
association. Incomparable directives needing one winner are refused, not
ordered by timestamps. There is no synthetic global commit DAG.

\paragraph{Partial fleet release.}
The protocol captures repository heads and tag refs, rechecks relevant inputs
before publication, and acquires fleet locks in a stable order. Native release
records and peer-link settlement admit their exact new heads; unrelated changes
require replanning. Peer links are settled before consumer publication so its
record carries the boundary evidence. Source tags remain authoritative after
failure: a failed control checkpoint requires explicit repair, not republishing
the successful source. Recording failures block dependent work while preserving
prior successes. This is forward recovery across separate ledgers, not an
atomic fleet transaction. These source-reviewed fleet mechanisms extend the
single-history model; the campaigns below do not evaluate them.

\subsection{Integration Limits}

A Go consumer resolving a new module version needs its tagged source published;
an upstream compiled output does not provide it. Local \code{go.work} or
\code{replace} can instead bind a checkout~\cite{gomod}. This boundary is
configuration-specific and unevaluated here.

Conforming integrations supply stable identities, declared dependencies and
availability, commands whose success means stage completion, and trustworthy
durable records. Uncertain publication requires safe repetition or
reconciliation. A specification can use validation/rendering as its build
and bundle publication as its publish stage. The model assumes an acyclic
graph and SemVer-compatible versions; arbitrary workflows need not conform.

\paragraph{Shared release logic across operating systems}
A game developer can use Windows with Linux dedicated-server CI, or Linux
with Windows packaging CI. Shared configuration holds the graph, lifecycle,
and portable commands; referenced platform files choose \code{sh} or PowerShell
and override platform-specific scripts. This reuses orchestration logic,
without translating shells or making outputs portable. Toolchains must exist
on each host. The evaluated conditional wrapper requires POSIX even when a
child selects PowerShell; a Windows caller without POSIX must select its
configuration directly. Native Windows equivalence and game-development
adoption are not measured.

\subsection{A Stable CI Entry Point}
The package graph drives named tests and releases through a generic CI job.
Adding a package changes configuration and scripts but may leave a provisioned
workflow unchanged. New runners, credentials, or system dependencies still
require provisioning. A configured gate must reject missing required tests
and block publication on failure.

Affected-task execution already appears in systems such as Nx~\cite{nx}. Here,
one package model supports testing and recoverable release execution across
artifact kinds. The two may use different windows and policies; the shared
graph avoids duplicate package lists without replacing CI policy.

\paragraph{Observed self-hosting case}
CI uses its checkout-built CLI for transitive test and build selection and
self-release. This single-project observation shows that the protocol can drive
its own repository; it is not adoption or causal productivity evidence.
